\section{Experiment}
%==================== SPEC (section-wide) ====================
%SPEC style: drop the literal "RQ" before \ref{rq:*} -- \ref already prints "RQ".
%DONE [refs 2026-06-29] swept paper-wide: removed all 11 "RQ"/"RQ~" prefixes before
%  \ref{rq:*}, and trimmed prose \autoref{sec:*} to <=1 per target (40 -> 11), keeping
%  only cross-chapter bridges; all tab/fig/eq refs untouched.
%SPEC reporting: surface the LLM backend + per-project results somewhere (manual, later).
%consolidate check / judge
%============================================================

We study four research questions:
\subsection{Research Questions}

% RQ1 Can the multi-stage LLM workflow outperform SOTA approach?
% \emph{\textbf{Motivation:}}

\begin{enumerate}[leftmargin=*, labelindent=0em, label=\textbf{RQ\arabic*}:,
ref=RQ\arabic*, series=researchQuestions] \item \label{rq:sota}
\emph{Does \approach recover trace links better than existing approaches?} 
\end{enumerate}
\emph{\textbf{Motivation:}}
Existing approaches either match component names lexically~\cite{keim_trace_2021} or classify with a single \ac{LLM} call~\cite{fuchs_whos_2025}. 
\approach builds runtime knowledge about the architecture and then runs focused linker modules with \evidenceBacked judges.
We want to know whether this extra structure helps in both link recovery tasks: doc-model, and its downstream task doc-code.

\begin{enumerate}[leftmargin=*, labelindent=0em, label=\textbf{RQ\arabic*}:, ref=RQ\arabic*, series=researchQuestions, resume]
    \item \label{rq:metrics} \emph{What do architecture-driven metrics tell us about the recovered links that the standard link-level metrics do not?}
\end{enumerate}
\emph{\textbf{Motivation:}}
The standard evaluation averages precision, recall, and \fone\ over all benchmark links.
This average is dominated by a few large components, hiding whether a system covers every sentence or drops small components.
We check whether the proposed architecture-driven metrics expose these properties and reorder the systems.

% RQ3 whats the contribution of validators in our approach?
% \emph{\textbf{Motivation:}}

\begin{enumerate}[leftmargin=*, labelindent=0em, label=\textbf{RQ\arabic*}:, ref=RQ\arabic*, series=researchQuestions, resume]
    \item \label{rq:validators} \emph{How much do the judges inside \approach contribute to the final result?}
\end{enumerate}
\emph{\textbf{Motivation:}}
% \acp{LLM} hallucinate: a linker can return a link that looks plausible but does not hold up under inspection.
To reduce hallucinated links, \approach{} wraps each link proposer in its corresponding \evidenceBacked judge. 
We measure how much each judge helps by calculating the confusion matrix.

% RQ4 ablation, whas the contribution of each sub-linker module?
% \emph{\textbf{Motivation:}}

\begin{enumerate}[leftmargin=*, labelindent=0em, label=\textbf{RQ\arabic*}:, ref=RQ\arabic*, series=researchQuestions, resume]
    \item \label{rq:ablation} \emph{What does each module of \approach{} contribute on its own?}
\end{enumerate}
\emph{\textbf{Motivation:}}
%%TODO, simplify
\approach{} splits the linking work by \emph{evidence form}---a whole name, one word of a name, or a referring expression---and gives each form the proposer that form admits and a judge of its own, over one shared knowledge module.
Whether that decomposition pays is a different question from whether the judges pay (\ref{rq:validators}): a judge can only remove links, whereas the form decides which links exist to remove.
A module-level ablation over the alias discovery module, and two linker modules.

\subsection{Experiment Design}
\label{sec:exp-design}

To answer the four research questions we run four families of experiments on the the \ardoco benchmark, the largest benchmark to the best of our knowledge~\cite{fuchs_establishing_2023}.  
It contains five Java open-source systems with hand-curated architecture docs, an architecture model, and a gold standard at two grains, document-to-component and document-to-code.
%DONE [dataset-stats] (2026-08-31) the "basic stat" this used to ask for is
%  \autoref{tab:gold_concentration}: per project, sentences in the SAD, components in the
%  model, gold-reachable components, and the link-concentration columns. Its own caption
%  says it doubles as the dataset table for this section, so no separate table is needed
%  and the space concern is moot. Generated by mini-inequality/motivation.py; do not
%  hand-edit table/gold_concentration.tex.
All baselines are run with their authors' released configurations.
% \emph{MediaStore}, \emph{TeaStore}, \emph{Teammates}, \emph{BigBlueButton}, and \emph{JabRef}.
% We use this benchmark because it is the standard evaluation set for this task~\cite{keim_recovering_2024,fuchs_lissa_2025,fuchs_whos_2025}.
For \ac{LLM} systems \approach{} and \Artemis{}, because they are non-deterministic, we run them three times and report the mean, so neither is credited with a lucky chance. 
The deterministic baselines are run once.
To avoid depending on one model, we run \approach{} on two language-model backends, GPT-5.6-terra and GPT-5.6-luna. 
% The body reports terra and the appendix mirrors every table on luna.

\subsubsection{Experiments for \ref{rq:sota}}
\label{sec:exp:rq1}
The two evaluations are:
\begin{enumerate}[leftmargin=*]
    \item \emph{Doc-model.} For every sentence in the architecture document, we use \approach to link it back to architectural components.
     We then compare recovered links with the doc-model gold standard links. 
    We compare \approach{} with two baselines: \emph{SWATTR}~\cite{keim_trace_2021}, a lexical string-similarity pipeline; and \Artemis{}~\cite{fuchs_whos_2025}, an \ac{LLM} entity-recognition tool for the same task.
    This is the task where \approach{} can directly show its strength.
    \item \emph{Doc-code.} Doc-code is the downstream task of doc-model: following how prior work~\cite{fuchs_whos_2025,keim_recovering_2024} were evaluated, \approach{} recovers the doc-model links first, then composes them into code links. 
    For every sentence we recover the code files it talks about and compare them with the link-level gold standard. 
    \approach{} reaches code files by transitive linking, the same composition used by \TransArc{} and Artemis\cite{fuchs_whos_2025}.
    It combines its doc-model links with the architecture model's component-to-file link recovery via ArCoTL~\cite{keim_recovering_2024}. 
    We compare \approach{} with two baselines: \TransArc{}~\cite{keim_recovering_2024}, the strongest heuristic-based approach; 
    and \Artemis{} composed through the same \TransArc{} file mapping, as an \ac{LLM} baseline at the doc-code grain.
\end{enumerate}
For each task we use the same metrics \TransArc{} and \Artemis{} use.
We report link level precision, recall, and \fone, \ftwo per project and averaged across projects.
% We report \ftwo\ beside \fone\ throughout, as previous studies argue both are valuable traceability tasks\cite{todo}.
The outcome of \ref{rq:sota} tells us whether \approach{} improves over existing tools on the link task it targets directly (doc-model) and on the downstream task that developers care about (doc-code).


\subsubsection{Experiments for \ref{rq:metrics}}
\label{sec:exp:rq2}
\ref{rq:sota} scores every system with the standard link-level metrics.
The pre-study (\autoref{sec:metric:prestudy}) shows these metrics are limited because of the benchmark's extreme inequality.
\ref{rq:metrics} re-runs the same comparisons under the metric suite of \autoref{sec:metric:suite}.

%DONE, \cmr is justified within the suite, so this reads as "apply the suite" and nothing
%  more. \cmr is the doc-model-only size-aware metric; it is not used on the file grain.
For each (system, project) cell on the doc-code task we compute the \fone-tail of the suite: reference link-level \fone\ and the two component-size aware metrics (worst-component and harmonic-mean per-component \fone).
On the doc-model task, we report the \cmrname{} (\cmr, \autoref{eq:cmr}), the weighted rate of abandoned components.
% We read the size-aware metrics beside link-level \fbeta: where a system scores well on the standard metric yet poorly on the size-aware ones, it has recovered the bulk of the links while leaving whole components untraced.

\subsubsection{Experiments for \ref{rq:validators}}
\label{sec:exp:rq3}
We build three ablations of \approach{} that remove judges, one per judge plus the union.
Running on the same benchmark, we compare each ablation against \emph{\fullVariant{}}, the full pipeline.

\begin{enumerate}[leftmargin=*]
\item \emph{\noNameValid{}}: skip the \nameValidator{}; accept every candidate of \linkerN{}, from both the \linkerB{} and the \linkerD{}, unjudged.

\item \emph{\noCitation{}}: skip the \corefValidator{}; accept all implicit-reference candidates without \corefValidator{}.

\item \emph{\noValidator{}}: skip both judges (the two above combined).
\end{enumerate}
Each judge logs a verdict per candidate it sees, so we read the ablations off the full pipeline's own decision log rather than re-running it three times.


% For each ablation we report (i) the change in macro \fone\ over the five projects against \fullVariant{}, 
% (ii) the true positives the validator preserves and the false positives it removes against the unvalidated linker output (the net validator effect), 
% and (iii) how that net effect breaks down per component, 
% to show whether the validator filters uniformly or only on certain components.

\subsubsection{Experiments for \ref{rq:ablation}}
\label{sec:exp:rq4}
%TODO, decide should we use both venn and table. 
Different from \autoref{rq:validators} on judge level ablation, we ablate \approach{} at the module, a higher level in this rq.
The workflow reaches three reference forms through two linkers over one knowledge module, so the ablation is by \emph{form}: the \linkerB{} and the \linkerD{} inside \linkerN{}, the \linkerC{}, and the alias table all three of them read.
Each link records the proposer that reached it, so the three forms can be separated after the fact even though two of them are judged together.
We measure their contributions with different methods.
For each form we measure that by set overlap.
This reports how many unique true positive links each form contributes that the others do not.
For the knowledge module we disable it, re-run and compare against the full workflow.
This isolates how much of the gain comes from the shared knowledge rather than from the linkers themselves.
\autoref{tab:rq4} reports the size-aware suite for the full workflow, for each form on its own, and for the no-knowledge variant.

%NOTE [dead block] (2026-08-31) this whole \subsection{Dataset} is SUPERSEDED and stays
%  commented out only as provenance: its live text was absorbed into the opening of
%  sec:exp-design above, and its dataset-stats TODO is answered by tab:gold_concentration
%  (see the %DONE [dataset-stats] note there). Nothing below this line is compiled. Delete
%  the block at the final pass if it is still here.
% \subsection{Dataset}
% \label{sec:dataset}

% We use the ARDoCo benchmark, the largest architectural traceability benchmark to the best of our knowledge~\cite{fuchs_establishing_2023}.
% It contains five Java open-source systems with hand-curated architecture docs, an architecture model, and a gold standard at two grains, document-to-component and document-to-code:
%  (superseded copy of the dataset-stats note -- answered by tab:gold_concentration)
% \emph{MediaStore}, \emph{TeaStore}, \emph{Teammates}, \emph{BigBlueButton}, and \emph{JabRef}.
% We use this benchmark because it is the standard evaluation set for this task~\cite{keim_recovering_2024,fuchs_lissa_2025,fuchs_whos_2025}.

 % lists, per project, the number of sentences in the architecture document and the number of components in the architecture model, alongside the link-concentration statistics analyzed there.
% \autoref{sec:metric:prestudy} discusses the two structural properties of this expansion: the per-project expansion factor and the long-tail per-component link distribution, that shape what the standard link-level \fone\ actually measures.

% For each project we use the architecture document as input~$D$, the architecture model as the component set~$C$, and the gold standard as the ground-truth links between sentences of~$D$ and components of~$C$ (or files belonging to components of~$C$, for the doc-code task).
